\section{A direct crossing and the exact joined arithmetic windows}\label{sec:windows}
The preceding character obstruction specifies what a successful crossing must
change. There is also a concrete positive crossing, and it extends to a sharp
joint window on an original factor packet. None of the following constructions
adds a factor that is absent from the source integer.

\begin{theorem}[The maximal divisor domain of a channel involution]\label{thm:quarter}
At an original source with $4\mid a$, put $V=\operatorname{Div}(a/4)$. The map
\[
 \mathcal T(u)=\frac a{4u}\quad(u\in V)                         \tag{W1}
\]
is an involution of $V$ and interchanges the tagged E and M hit sets on $V$.
Among all original divisors $u\mid a^2$, this is its complete integral return
domain. Retaining the channel tag, its inverse is the same formula. Its only
possible numerical fixed points satisfy $a=4u^2$; they are not identified across
the two channel labels.
\end{theorem}
\begin{proof}
$\mathcal T(u)$ is an integer if and only if $4u\mid a$. In that case it and $u$
divide $a/4$, and the second application returns $u$. Every such divisor is a
unit modulo $R$. Thus
\[
 R\mid4u+1\iff R\mid a(4u+1)=4u(\mathcal T(u)+a)
 \iff R\mid\mathcal T(u)+a.
\]
The character ratio is $\chi_p(\mathcal T(u)/u)=\chi_p(a)$ because
$\mathcal T(u)/u=a/(4u^2)$. Thus this is the actual non-square crossing
required on a negative source, rather than a forbidden square-ratio transport.
This also proves the reverse channel assertion. The fixed-point equation is
$4u^2=a$. In particular negative-character sources have no numerical fixed
point. The positive-source example $p=61,a=16,R=3,u=2$ is a fixed point, with
two different tagged denominator returns.
\end{proof}

Write $a=2^e A$, $A$ odd, and suppose $e\ge1$. Retain only the specified packet
\[
 \mathscr D=\{2^fd:0\le f\le2e,\ d\mid A\}\subseteq\operatorname{Div}(a^2).
                                                               \tag{W2}
\]
The omitted original divisors, whose odd parts divide $A^2$ but not $A$, remain
outside this packet, not outside the complete atlas. For each $w\mid A$ define
\[
 I_E=[-2e-2,-2]_{\mathbb Z},\quad I_M=[-e,e]_{\mathbb Z},\qquad
 2^j\equiv-w\pmod R.                                      \tag{W3}
\]
\begin{theorem}[Joined-window coefficient theorem]\label{thm:windows}
The maps
\[
 (w,j,E)\longmapsto u_E=2^{-j-2}w,\qquad
 (w,j,M)\longmapsto u_M=2^{e+j}A/w                         \tag{W4}
\]
are bijections from the respective solutions of (W3) in $I_E,I_M$ onto the
original tagged packet hits. Their inverses read $f=v_2(u)$, the odd part $d$,
and return $(w,j)=(d,-f-2)$ for E and $(A/d,f-e)$ for M.
The union of the two integer windows is the interval
\[
 [-2e-2,e]_{\mathbb Z},\qquad |I_E\cup I_M|=3e+3.          \tag{W5}
\]
On their intersection $[-e,-2]$ the two divisors have product $a/4$, giving
exactly the crossing (W1). For $e=1$ that intersection is empty.

Let $d_R=\operatorname{ord}_R(2)$ and
$A_- =\#\{w\mid A:-w\in\langle2\rangle_R\}$. If $3e+3\ge d_R$, the packet
is jointly occupied if and only if $A_->0$. Moreover put
\[
 2e+1=f_0d_R+\rho,\quad0\le\rho<d_R,\quad
 \Delta=(e+2)\bmod d_R,\quad\delta=\min(\Delta,d_R-\Delta).
\]
Then its exact joint count $C_{\rm packet}$ obeys
\[
 A_-\bigl(2f_0+\mathbf1_{\rho\ge d_R-\delta}\bigr)
 \le C_{\rm packet}\le
 A_-\bigl(2f_0+\mathbf1_{\rho>0}+\mathbf1_{\rho>\delta}\bigr). \tag{W6}
\]
The bracketed constants are the minimum and maximum of the paired cyclic
count at a target class. In particular every cyclic target is covered exactly
when $3e+3\ge d_R$; this is an exact interval-cover threshold, not an assertion
that every abstract target is realized by a prime source.
\end{theorem}
\begin{proof}
For E, $4u+1=0$ is equivalent to $2^{-f-2}=-d$. For M, divide
$2^fd+2^eA=0$ by the unit $2^ed$ to obtain $2^{f-e}=-A/d$.
This proves (W3)--(W4) and their inverses. Their integer intervals have no gap
when $e\ge1$. On the overlap the displayed product is $2^{e-2}A=a/4$.

For the count, remove $f_0$ complete periods from each window. The two remaining
cyclic intervals have length $\rho$ and displacement $\Delta$. They overlap
precisely when $\rho>\delta$, and cover the whole cycle precisely when
$\rho\ge d_R-\delta$. Thus the minimum of their indicator sum is zero or one
as in (W6), and its maximum is zero, one or two as stated. Sum this exact bound
over the $A_-$ actual divisors; a divisor outside the subgroup contributes zero.
For coverage, (W5) contains a complete period precisely at its stated length.
\end{proof}
The two label occurrences in an overlap are not one state. For $e=0$ the
windows are the singleton sets $\{-2\}$ and $\{0\}$; (W5)--(W6) are not used.
The same prime's upper-half sources remain in the data; their M packet is empty
by the already proved denominator-size argument.

\section{A complete classification when the nonresidue factor is simple}
The joined windows admit a complete rather than subpacket interpretation on a
specified original factor class. This separates an actual odd-component hole
from merely insufficient bounded powers.

\begin{theorem}[One simple nonresidue factor]\label{thm:simpleNR}
Suppose $a=Br$, $r$ is prime, $r\nmid B$, $(r/p)=-1$, and every prime factor of
$B$ is a residue modulo $p$. Let
\[
 \mu_B(g)=\#\{v\mid B^2:vB^{-1}=g\pmod R\}.
\]
Then the \emph{complete} original channel counts are
\[
 \boxed{|E_a|=\mu_B(-p^{-1}),\qquad |M_a|/2=\mu_B(-r).}  \tag{W7}
\]
The exact E fibre is $u=rv$, with $vB^{-1}=-p^{-1}$. The first M orientation
is $u=v$, with $vB^{-1}=-r$, and its paired orientation is $a^2/v$.
The complements are distinct and exhaust M. All original exponent vectors of
$B$ and their fine residue multiplicities are retained.
\end{theorem}
\begin{proof}
The original source has character $-1$. The channel-character theorem forces
an odd exponent of $r$ for E and an even exponent for M. The only permitted
exponents are $0,1,2$. At exponent one, $u/a=v/B$, giving the E formula.
At exponent zero, $u/a=v/(Br)$, so M requires $v/B=-r$. Complementation sends
this orientation to exponent two and is a bijection, proving the full count.
\end{proof}

In particular let $B=\ell^e$, with $\ell$ prime and $(\ell/p)=1$. Put
$H=\langle\ell\rangle_R$, $d_R=|H|$. If $4\in H$, retain
$c=\log_\ell 4\in[0,d_R)$; if $-r\in H$, retain $j=\log_\ell(-r)$.
The complete counts become
\[
 \boxed{\begin{aligned}
 |E_a|&=\#\{j'\in[-2e-c,-c]:j'\equiv j\pmod{d_R}\},\\
 |M_a|/2&=\#\{j'\in[-e,e]:j'\equiv j\pmod{d_R}\}.
 \end{aligned}}                                           \tag{W8}
\]
If $-r\notin H$, both vanish. The E map is $u=\ell^{-j'-c}r$, and the first
M map is $u=\ell^{e+j'}$, with the complementary M word retained. Each window
has length $2e+1$, hence
\[
 \boxed{\bigl||E_a|-|M_a|/2\bigr|\le1.}                  \tag{W9}
\]
If $4\notin H$, M requires $-r\in H$, while E requires $-r\in4^{-1}H$.
Those are disjoint cosets, so at most one channel can occur. No nonexistent
logarithm $c$ is assigned in that case.

For $\ell=2$, one may use $c=2$ as an integer lift regardless of the order,
and Theorem~\ref{thm:windows} is an exact classification of $|E|+|M|/2$ on
this factor class, not just a subpacket. The full M count is twice its first
orientation count. At fixed $R$ and fixed odd factor $A$, increasing $e$ by
$d_R$ adds two occurrences per target to each labelled packet in (W3). This is
an all-exponent recurrence with $p_e=4\cdot2^e A-R$ \emph{varying}; primality
and the original source range must be checked for each parameter.

The cyclic threshold has an actual sharp boundary example:
\[
 p=23209,\quad R=127,\quad a=2\cdot2917,\quad
 d_R=7,\quad-r\equiv2^2\pmod{127}.
\]
The E window occupies $3,4,5$, the M window $6,0,1$, and class 2 is absent.
Both complete channels are empty one class below $3e+3\ge d_R$.
At the hard primes $43801,8521,12601$, with residual 31 and respectively
$a=2\cdot5479,2\cdot1069,2\cdot1579$, the complete pairs
$(|E|,|M|/2)$ are $(1,0),(0,1),(1,1)$. The one-unit difference in (W9)
cannot be removed. These primes and every factor are checked by trial division.

\section{Six expanded vertices can still have empty complete square-source channels}
For a prime vertex $q$, write $\sigma_q=1$ at $q\equiv3\pmod4$ and
$\sigma_q=3$ at $q\equiv1\pmod4$. Its square source has
$R=\sigma_q qt^2$, $a=(p+R)/4$, and positive odd $t$ with $R<3p$.
Turn 4 forces six distinct reachable vertices, not a selector. Here is an exact
local obstruction at that cardinality, beyond the preceding two-vertex example.

\begin{theorem}[A complete six-vertex local failure certificate]\label{thm:six}
At the hard prime $p=315361\equiv361\pmod{840}$, the canonical graph has the
closed six-cycle
\[
 27847\to42901\to13877\to2789\to20233\to18803\to27847.    \tag{W10}
\]
Every vertex has exactly its displayed outgoing nonresidue prime, to exponent
one. At all \emph{17} valid square sources of these six vertices, both original
E/M channels are empty. The complete source consists of 663 original divisors.
This is not a closed set in the expanded square graph and is not an ES
counterexample.
\end{theorem}
\begin{proof}
The six complete canonical factorizations are
\[
 2\cdot42901,\quad2^3\cdot13877,\quad2^5\cdot2789,\quad
 2^2\cdot20233,\quad5\cdot18803,\quad3\cdot27847.
\]
Exact trial division proves $p$ and all factor primes. Their quadratic symbols
prove the stated complete successor set. The finite t-ranges, in that order,
are $\{1,3,5\}$, $\{1\}$, $\{1,3\}$, $\{1,3,5,7,9\}$, $\{1,3\}$, and $\{1,3,5,7\}$.
They follow from the strict inequality $\sigma_q qt^2<3p$. The certificate lists
every factorization and every bounded divisor and checks both gates directly.
Every source in fact has a unique simple nonresidue factor, so (W7) independently
reduces the two full gates to the two exact positive-factor coefficients.
Both reductions are checked against the complete divisor list.

For the canonical source $q=42901$, the pure-$2$ subgroup modulo $128703$
has order 14300 and contains neither required target. Several other sources
have their target in the generated subgroup but outside their actual bounded
packet. These are two different certified reasons for vanishing.
\end{proof}
An actual exit and successful next source are retained. At $q=20233,t=3$,
\[
 A=215413=11\cdot19583
\]
has the nonresidue factor 11. The canonical source at that prime is
$A_{11}=78843=3\cdot41\cdot641$ and has an exterior hit $u=41$ as well as
four middle states. Its E normalization and ordered triple are in the certificate.
This is an actual path, not a proof that every newly produced factor reaches a hit.
No leastness or complete-square-graph failure is asserted for (W10).

The finite certificate is independent of the discovery search that suggested
this cycle. The proof only needs the displayed prime, its six complete factor
identities, the exact valid source ranges, and all 663 divisor vectors. The
independent checker uses direct integer divisors and literal square-residue sets,
not the first implementation's signed-box recurrence or Euler-symbol routine.
